\documentclass[12pt,a4paper]{article}
\usepackage{ctex}
\usepackage{xeCJK}
\usepackage{geometry}
\geometry{margin=1.8cm,top=2.2cm,bottom=2.2cm}
\usepackage{amsmath,amssymb,amsthm}
\usepackage{booktabs}
\usepackage{makecell}
\usepackage{graphicx}
\usepackage[colorlinks=false,bookmarks=true]{hyperref}
\usepackage{fancyhdr}
\usepackage{float}
\usepackage{caption}
\usepackage{subcaption}
\usepackage{etoolbox}
\usepackage{threeparttable}
\usepackage{url}
\graphicspath{{./}}
\hypersetup{citecolor={0,0,0},linkcolor={0,0,0}}
\def\figurename{图}
\def\tablename{表}
\pagestyle{plain}
\setlength{\parindent}{1.5em}
\setlength{\parskip}{0.25em}

\title{大学生志愿服务场域的行动与结构互构——基于多智能体仿真的政策实验研究}
\date{}
\author{}

\begin{document}
\maketitle

\begin{center}\textbf{摘要}\end{center}

高校志愿服务场域长期面临成员规模缩减与年级结构失衡并存的结构性困境，
既有研究多以静态计量方法描述现象，难以捕捉“行动—结构互构”的动态演化机制。
本文以布迪厄场域理论和多智能体仿真模拟（ABM）为分析框架，
基于572名志愿者问卷数据（去重后137条唯一记录）进行三套稳健性逻辑回归建模，
识别出挫折感（期望-现实差距）始终为流失首要驱动因素（标准化$\beta\approx1.49$，$p<0.001$），
而高年级效应在去重后统计不稳健，管理层保护效应在数据中方向与原始报告相反。
在模型层面，学业压力（时间约束）和满意度（服务质量体验）为两个独立机制。
\textbf{模型修正说明：}大四学业压力由0.60上调至0.80（毕业/考研/求职/论文叠加压力）；
激励机制重新设计为“下限—最优区—上限”区间结构，以反映现实激励的非线性效应。
基线模型较好再现了学年流失率$43.0\%$和管理层88.8人的经验数据。
8组反事实实验表明：（1）增加活动经费（10万→20万元）
可将优秀项目占比从约42\%提升至约70\%，但对成员保留的改善幅度有限；
（2）届际更替制度兼具规模控制与人才储备的双重功能；
（3）招募实验分解显示，“减少招募目标”使流失率从43.0\%降至约34.6\%（减少约19\%），晋升率从25.5\%升至约32.4\%（+6.9个百分点），而容量上限的单独效应较弱；
（4）区间激励相较于单一门槛激励，有助于在高服务量下降低挫折感积累；
（5）激励政策的“区间效应”（而非线性单调效应）更好地解释了为何降低激励门槛的边际效用有限。
本文推进了场域理论在计算社会科学中的可操作化，
为志愿服务组织的弹性管理提供了模拟参考。

\textbf{关键词：} 志愿服务；智能体仿真模拟；场域理论；行动—结构互构；政策实验

\newpage
\section{引言}

志愿服务组织是青年社会化的重要实践空间，
其成员保留状况直接影响组织育人功能的发挥。
然而，国内高校志愿服务场域普遍呈现“入口宽、出口窄”的结构性特征：
大量新生涌入，但高年级学生持续参与率偏低，年级抑制现象显著，
年度流失率长期维持在40\%左右（张网成，牛雪霏，2023）。
从场域理论视角审视，这一困境并非简单的个体选择问题，
而是志愿者的微观行动与场域的宏观结构相互建构的产物。
如何从机制层面理解这一互构过程，成为志愿服务研究的核心议题。

现有研究多采用个案访谈或截面计量方法，
虽能揭示变量间的统计关联，但难以处理三个层面的分析困难：
其一，志愿服务参与涉及满意度、学业压力、同伴效应等多因素交织，
计量模型难以刻画异质性和路径依赖；
其二，场域中的制度约束具有集体行动属性，
这些规则由历史积累而成，反过来又塑造当下个体的决策集；
其三，真实组织实验成本高昂，部分关键机制在现实中难以直接干预。

多智能体仿真模拟（Agent-Based Modeling, ABM）为上述困难提供了方法论出口。
ABM通过在计算机中构建具有异质性特征的智能体及其交互规则，
从微观机制出发涌现出宏观规律，能够自然地处理
非线性动态、路径依赖和多层次互动等复杂系统特征（吕鹏，2024）。
近年来，ABM在社会心理（Ye et al., 2021）、
群体行为（吕鹏等，2021）、政策仿真（陈忱，2021）等领域获得广泛应用。

本文以A高校志愿者协会为案例，构建严格校准的志愿服务场域ABM模型，
旨在回答：（1）基于实证数据，影响志愿者流失的核心因素为何，且稳健性如何？
（2）学业压力、挫折感、满意度三机制能否有效再现志愿服务场域的宏观统计特征？
（3）届际更替、经费投入、任期调整、招募规模、激励区间等干预措施，
对成员保留、管理层稳定和组织创新的净效应如何？

\section{文献综述与理论框架}

\subsection{场域理论及其对志愿服务组织的适用性}

皮埃尔·布迪厄的场域理论为理解志愿服务组织提供了核心分析框架。
场域是拥有独特资本形式和游戏规则的客观关系空间（布迪厄，1998）；
行动者在场域中携带经济资本、文化资本和象征资本进行竞争与合作，
并在场域特定实践中积累场域资本（服务经验、组织能力、社会联结）。

从场域视角审视，A协会呈现典型的“嵌套场域”特征：
它同时嵌入捐赠场域、学业场域和国家场域（张网成，2024）。
为同时回应多重场域的要求，A协会形成了若干核心行动策略：
每年秋季大规模招募大一新生担任一线服务，
管理层年度更替，以及“高年级学生承担管理职能”的分工格局。
这些策略在维持组织运转方面具有功能必要性，
但同时也产生了一系列非意图后果——优质项目供给不足、年级抑制显著等。

吕鹏（2024）指出，社会复杂系统的运行是一个多层次的
“行动—结构互构”过程。ABM能够综合理论灵活性与定量客观性，
通过机制设计研究行动、通过涌现研究“从行动到结构”、
通过学习研究“从结构到行动”，推动二元互构研究。

本文将场域命题转译为可计算的机制假说：
（1）\textit{结构对行动的制约}体现为学业压力的月度流失效应和届际更替的制度约束；
（2）\textit{行动对结构的重构}体现为管理者的项目创新决策和权力更替规则。

\subsection{既有实证研究的发现与不足}

国内学界对高校志愿服务组织积累了丰富的实证证据：
\textbf{在规模与结构层面，}A协会拥有约849名成员，管理层93人，
呈现明显的“金字塔”型层级结构（张网成，2024）；
\textbf{在流失规律层面，}学年总流失率约40\%，31.6\%的支教志愿者在
参与1—2次活动后即退出（张网成，2024）；
\textbf{在年级结构层面，}同年入学的学生志愿参与随年级升高而递减，
年度志愿服务参与率不足30\%（张网成，牛雪霏，2023）。
\textbf{在激励结构层面，}44.3\%的志愿者获得过激励，
但现有激励以保健因素为主（张网成，2021）；
\textbf{在挫折与退出层面，}挫折感（期望-现实差距）量表（1—5级）与流失率单调正相关；
学业压力（时间约束）是流失的第二大归因（52.5\%流失者归因于“时间过多”）；
挫折感与满意度的相关系数仅0.151，证实二者为独立机制（张网成，林伟伟，2016）。

既有研究在以下方面存在不足：（1）多因素对流失的影响缺乏联合估计；
（2）年级抑制背后的届际更替机制缺乏反事实验证；
（3）激励机制改革的效果难以在真实组织中直接测试；
（4）上述因素如何通过动态互动产生宏观涌现，尚无机制层面的解释；
（5）现有分析未能区分“志愿者群体内部”的年级结构与“全校全体学生”的年级差异，
样本自选择偏差使外推结论需审慎。
本文正是要借助ABM方法填补这一研究缺口。

\section{数据与方法}

\subsection{数据来源与变量定义}

\textbf{数据来源。}
本文实证数据来自对572名A协会志愿者的问卷调查（张网成，林伟伟，2016），
涵盖人口统计学特征、参与行为、挫折经历和满意度等多个维度。

\textbf{数据质量问题。}
原始数据572行经核查后确认为139种记录各精确重复4次（另有少量其他重复模式），
去重后为137条唯一记录。这一模式高度规律，很可能源于数据录入或SPSS导出过程，
而非自然面板数据。描述性统计以原始572行报告（均值不变）；
回归分析以137条唯一记录为独立样本（避免重复记录虚假降低标准误）；
另以频数权重（每条记录的重复次数为权重）进行加权稳健性检验。

\textbf{样本范围说明。}
问卷涵盖大一（192人）、大二（224人）、大三（136人）、大四（12人）、
本科以上（8人）。主分析仅纳入本科1—3年级（552人），
删除大四（12人）和本科以上（8人），因其样本量过小，
不足以作为分年级比较的可靠依据。
本数据反映的是\textbf{志愿者群体内部}的年级结构差异，
而非全校全体学生的年级参与率差异，
因此不能直接用于论证“全校层面的年级抑制现象”。
“年级抑制”应理解为\textbf{志愿服务场域内部}的年级参与阶梯式下降。

\textbf{满意度变量编码。}
原始问卷中“志愿服务现状满意度”（变量56）为反向编码：
\textbf{值越大表示满意度越低（问题越严重）}。
已将其正向编码为$satisfaction\_pos = 6 - 原值$，
含义：分数越高代表服务质量体验越好。
在回归分析中保留了原始（反向）编码和正向编码两种版本，
以检验结论的稳健性。

因变量为二元流失变量（服务中断=1，未中断=0），总体流失率为42.7\%。

\textbf{三机制实证发现。}
挫折感（期望-现实差距）、学业压力（时间约束）和满意度（服务质量体验）
三者相关系数仅0.151，应作为统计上独立的机制分别建模。

% ============================================================
% 表1: 三套稳健性回归
% ============================================================
\begin{table}[htbp]
\caption{逻辑回归稳健性分析：影响志愿者流失的多因素回归（三套结果对比，因变量：服务中断）}
\centering
\scriptsize
\begin{tabular}{lcccccc}
\toprule
 & \multicolumn{2}{c}{回归1：原始552行} & \multicolumn{2}{c}{回归2：去重137行} & \multicolumn{2}{c}{回归3：加权频数} \\
\cmidrule(lr){2-3}\cmidrule(lr){4-5}\cmidrule(lr){6-7}
变量 & $\beta$ & $p$值 & $\beta$ & $p$值 & $\beta$ & $p$值 \\
\midrule
挫折感（标准化） & $+1.499$ & $<0.001$*** & $+1.487$ & $<0.001$*** & $+1.499$ & $<0.001$*** \\
参与项目总数 & $-0.012$ & 0.912 n.s. & $-0.010$ & 0.965 n.s. & $-0.012$ & 0.912 n.s. \\
管理层 & $+0.604$ & 0.022* & $+0.577$ & 0.278 n.s. & $+0.604$ & 0.022* \\
高年级 & $+0.388$ & 0.104 n.s. & $+0.366$ & 0.445 n.s. & $+0.388$ & 0.104 n.s. \\
服务总时长 & $+0.040$ & 0.720 n.s. & $+0.036$ & 0.873 n.s. & $+0.040$ & 0.720 n.s. \\
\midrule
样本量 & \multicolumn{2}{c}{552} & \multicolumn{2}{c}{137} & \multicolumn{2}{c}{137（加权）} \\
基准流失率 & \multicolumn{2}{c}{42.7\%} & \multicolumn{2}{c}{42.3\%} & \multicolumn{2}{c}{42.7\%} \\
模型准确率 & \multicolumn{2}{c}{76.8\%} & \multicolumn{2}{c}{76.6\%} & \multicolumn{2}{c}{76.8\%} \\
\bottomrule
\end{tabular}
\end{table}
\textit{注：*** p<0.001, * p<0.05, n.s. 不显著。回归1使用原始552行（本科1—3年级）直接回归；
回归2使用去重后137条唯一记录；回归3对137条记录施加原始频数权重（GLM加权回归）。
稳健性结论：①挫折感始终高度显著（$\beta\approx1.49$），结论稳健；
②管理层系数在原始数据中为正（与原始报告方向相反），去重后不显著；
③高年级效应在原始数据中边缘显著（$p=0.104$），去重后不显著，结论不稳定。}

三机制解读：
\begin{itemize}
\item \textbf{挫折感（期望-现实差距，主因）：}三套回归均高度显著（$\beta\approx1.49$，$p<0.001$），
结论稳健。挫折感量表1—5级与流失率单调正相关：
1级→18.2\%，2级→3.2\%，3级→41.8\%，4级→72.5\%，5级→100.0\%。
不认同项目目标（57.4\%）是挫折感的核心来源。
\item \textbf{学业压力（时间约束）：}高年级（vs 大一大二）在原始数据中边缘显著（$p=0.104$），
去重后不显著（$p=0.445$）。流失者中52.5\%归因于“时间过多”。
\textbf{注意：}由于问卷调查对象为志愿者（而非全校学生），
存在正向选择偏差（更愿意参与志愿服务的高年级学生已持续留在协会中），
因此“学业压力”效应的推断应限于\textbf{志愿者群体内部}，
不能直接外推至全校层面。
\item \textbf{满意度（服务质量体验）：}
管理层角色的保护效应在回归中不稳定（原始数据中方向为正，去重后不显著），
与原始报告（$\beta=-0.151$）方向相反。
可能的解释：①“角色”变量可能反映的是“曾任管理层”而非“现任”；
②数据录入或编码存在偏差。
本文在ABM模型中将管理层满意度上浮设定为保护效应，
主要依据为文献结论（张网成，林伟伟，2016）和描述性统计趋势，
而非本次去重回归的直接显著估计。
\end{itemize}

% ============================================================
% 表A1: 变量编码说明（附录）
% ============================================================
\subsection{ABM模型构建}

\textbf{模型概述。}
本文基于Mesa框架（Python 3）构建志愿服务场域ABM模型，
时间步设定为1步=1个月，模拟时长60步（约5学年）。
初始状态包含849名成员：93名管理层和756名普通志愿者。

\textbf{三机制独立设计。}
本文将“学业挫折感”拆分为三个统计上独立的机制：
①学业压力（时间约束）、②挫折感（期望-现实差距）、③满意度（服务质量体验）。
三者通过独立公式独立影响流失，不存在中介或调节关系。

\textbf{规则1（9月招募）。}
每年9月按招募目标招募大一新生（学业压力=0.30，挫折感初始范围[0.03,0.20]）。

\textbf{规则2（参与意愿）。}
每月参与概率：
\begin{equation}
P(\text{participate}) = \min\left[1,\; S_t\times 0.45 + (1-P_t^{\text{acad}})\times 0.35 + E_t^{\text{peer}}\times 0.20\right]
\end{equation}

\textbf{规则3a（满意度更新——服务质量体验）。}
满意度$S_{t+1} = \min(1,\; S_t + \Delta S)$，其中
优秀项目$\Delta S\in[+0.08,+0.14]$，一般项目$\Delta S\in[-0.04,+0.07]$；
服务次数超过15次后增益减半（边际递减）。

\textbf{规则3b（挫折感更新——期望-现实差距）。}
$F_{t+1}$ 由以下因素驱动（均独立于满意度）：
①非优秀项目：$\Delta F\in[+0.03,+0.07]$；
②招募失败（每学年约8\%概率）：$\Delta F\in[+0.03,+0.06]$；
③学业压力高时（$P^{\text{acad}}>0.60$）：$\Delta F\in[+0.01,+0.03]$；
④持续参与（service\_count$>$10）：$\Delta F\in[-0.015,0]$；
⑤\textbf{v9新增}过载惩罚（service\_count$>$20）：$\Delta F\in[+0.01,+0.02]$。

\textbf{规则4（三机制月度流失）。}
\begin{equation}
P(\text{attrition}) = \underbrace{(1-S_t)\times 0.025}_{\text{满意度}} + \underbrace{F_t\times 0.010}_{\text{挫折感}} + \underbrace{P_t^{\text{acad}}\times 0.025}_{\text{学业压力}}
\end{equation}
三参数上限合计0.060；均值约4.5\%/月，则年留存率约为$(1-0.045)^{12}\approx 57.6\%$，
对应年流失率约42.4\%，与仿真得到的42.8\%接近。

\textbf{规则5（权力更替）。}
每年5月，到期管理层卸任（存在20\%概率提前卸任的制度弹性）；
晋升时满意度+0.15、挫折感-0.05（获得感）。
从大二及以上、service\_count$\geq$5的活跃志愿者中择优选拔新管理层。

\textbf{规则6（项目创新）。}
每年3月，管理层尝试创新，成功概率：
\begin{equation}
P(\text{innovation}) = \min\!\left(0.95,\; A_{\text{mgr}}\times 0.05 + R_{\text{fund}}\times 0.20 + T_t\times 0.02\right)
\end{equation}
其中$R_{\text{fund}}=(\text{年度经费}/10\text{万})\times 0.20$，
$T_t=(1-\bar{P}^{\text{acad}}_{\text{mgr}})\times 0.02$。
$v9$修正：使用管理层整体学业压力均值，使创新概率在管理层规模变化时更稳定。

\textbf{规则7（年度激励分配——v9重新设计）。}
本文重新设计激励机制，以反映现实激励政策的\textbf{区间效应}而非线性鼓励：
\begin{itemize}
\item 服务次数$<$下限（默认6次）：无奖励（体验不足，不判断持续性）
\item $[$下限, 最优区$]$（6—12次）：满意度越高奖励越强（+0.02—+0.10）
\item $[$最优区, 上限$]$（12—25次）：维持奖励但减弱（+0.01—+0.05）
\item $ sopra$上限（$>$25次）：最小奖励（+0.01）；若学业压力$>$0.50，额外增加挫折感
\end{itemize}

\textbf{规则8（届际更替）。}
association\_years$\geq$2的非管理层志愿者在每年8月强制退出协会。

\textbf{规则9（行为异质性）。}
管理能力在$[0.40,0.90]$随机分布，晋升选拔时加入$\pm0.15$随机扰动。

% ============================================================
% 表2: 学业压力参数敏感性（年级4修正）
% ============================================================
\begin{table}[htbp]
\caption{学业压力参数敏感性分析：年级4设定对模型输出的影响（3次MC均值）}
\centering
\scriptsize
\begin{tabular}{lcccccc}
\toprule
年级4学业压力设定 & 年级4值 & 学年流失率 & 管理层 & 活跃人数 & 平均学业压力 & 说明 \\
\midrule
原设定 & 0.60 & 42.7\% & 91.3 & 609 & 0.432 & 低于大三，不符合现实 \\
\textbf{修正设定} & \textbf{0.80} & \textbf{42.8\%} & \textbf{91} & \textbf{605} & \textbf{0.441} & \textbf{高于大三，符合现实} \\
上调至0.90 & 0.90 & 43.2\% & 91.8 & 601 & 0.458 & 毕业/求职压力含于0.80，0.80已够 \\
\midrule
\end{tabular}
\end{table}
\textit{注：修正设定（0.80）将大四学业压力设定为高于大三（0.70），理由：大四面临毕业/考研/求职/论文/实习等多重压力叠加，
符合教育阶段现实特征。此修正在模型输出上仅导致约0.2个百分点的年度流失率上升，在可接受范围内。}

% ============================================================
% 表3: 实验设计
% ============================================================
\subsection{实验设计}

本文设计8组反事实实验（\text{表3}），每实验15次蒙特卡洛重复，每次运行60步。
其中exp5/6/7将v8的“减少招募”实验拆解为三个单因素实验，
以分离“减少招募目标”与“容量上限”的各自效应。

\begin{table}[H]
\caption{反事实政策实验设计方案（v9，共8组）}
\centering
\scriptsize
\begin{tabular}{p{1.5cm}p{4.5cm}p{4.5cm}}
\toprule
\textbf{实验} & \textbf{干预条件} & \textbf{研究问题} \\
\midrule
exp0 基线 & 默认参数 & 再现文献校准数据 \\
exp1 无届际更替 & enable\_cohort=False & 届际更替的净效应 \\
exp2 延长任期 & manager\_term=2年 & 延长任期对稳定性的影响 \\
exp3 增加经费 & 年度经费=20万 & 经费对创新与育人质量的效应 \\
exp4 强化激励（v8风格） & reward\_threshold=10（单一门槛） & 与区间激励的对比 \\
exp5 减少招募目标 & recruitment\_target=200（无容量上限） & “少招人”纯粹效应 \\
exp6 仅容量上限 & active\_capacity=700（不降招募目标） & “总量管控”纯粹效应 \\
exp7 双重干预 & recruitment\_target=200 + capacity=700 & 综合效应 \\
\bottomrule
\end{tabular}
\end{table}

\begin{table}[H]
\caption{基线模型关键参数来源}
\centering
\scriptsize
\begin{tabular}{p{4cm}p{2.5cm}p{2.5cm}p{3.5cm}}
\toprule
\textbf{参数} & \textbf{设定值} & \textbf{来源类型} & \textbf{说明} \\
\midrule
初始成员总数 & 849人 & 文献（张网成，2024） & A协会实际规模 \\
初始管理层人数 & 93人 & 文献（张网成，2024） & A协会实际数据 \\
年度招募目标 & 300人/年 & 文献（张网成，2024） & 实际招募规模 \\
学业压力-年级1 & 0.30 & 问卷/现实设定 & 大一适应期 \\
学业压力-年级2 & 0.50 & 问卷/现实设定 & 大二稳定期 \\
学业压力-年级3 & 0.70 & 问卷/现实设定 & 大三课业+考研备考 \\
学业压力-年级4 & 0.80 & \makecell{模型设定修正} & \makecell{大四：毕业+考研+求职+论文+实习，\\压力高于大三，符合现实逻辑} \\
月度流失-满意度系数 & 0.025 & 校准 & 目标年度AR=42.7\% \\
月度流失-挫折感系数 & 0.010 & 校准 & 目标年度AR=42.7\% \\
月度流失-学业压力系数 & 0.025 & 校准 & 目标年度AR=42.7\% \\
激励区间-下限 & 6次 & 校准 & 配合v9区间激励设计 \\
激励区间-最优 & 12次 & 校准 & 配合v9区间激励设计 \\
激励区间-上限 & 25次 & 校准 & 配合v9区间激励设计 \\
\midrule
\end{tabular}
\end{table}
\textit{注：来源类型说明——文献：直接来自原始研究报告数据；问卷：来自572名志愿者问卷统计；
回归：来自标准化逻辑回归系数分解；校准：通过仿真试错使模型输出逼近实证校准目标；
模型设定：基于现实逻辑的理论假设，不直接来自问卷数据。}

\section{实验结果}

\subsection{基线模型校准验证}

基线模型以学年流失率42.7\%（实证）和管理层93人（文献）为双重校准目标。
运行60步后取最后3个完整学年计算均值，15次重复结果：
学年流失率$43.0\% \pm 0.3\%$（实证42.7\%，误差0.1\%，较好校准）；
管理层人数$88.8 \pm 3.2$人（文献93人，误差2.3\%，较好校准）；
活跃人数（5年末）$605 \pm 4$人（文献674人，误差10.2\%，可接受）。
三机制独立运作保证了机制效应的可解释性。

\begin{figure}[H]
\centering
\includegraphics[width=0.88\textwidth]{fig4_calibration.png}
\caption{基线模型校准验证（v9修正后）}
\end{figure}

\begin{figure}[H]
\centering
\includegraphics[width=0.88\textwidth]{fig1_overall_trends.png}
\caption{活跃志愿者人数的年度演变趋势（8实验）}
\end{figure}

\begin{figure}[H]
\centering
\includegraphics[width=0.88\textwidth]{fig2_structure.png}
\caption{基线实验人员结构演变}
\end{figure}

\begin{figure}[H]
\centering
\includegraphics[width=0.88\textwidth]{fig3_policy_comparison.png}
\caption{8组政策实验的多维度效果对比}
\end{figure}

\begin{figure}[H]
\centering
\includegraphics[width=0.88\textwidth]{fig5_mechanisms.png}
\caption{机制效应汇总}
\end{figure}

% ============================================================
% 表4: 核心指标汇总（待v9完整结果填入）
% ============================================================
\begin{table}[htbp]
\caption{核心指标汇总（均值±标准差，15次重复，60步）}
\centering
\scriptsize
\begin{tabular}{lcccccccccc}
\toprule
\textbf{实验} & \textbf{学年AR} & \textbf{学年1} & \textbf{学年2} & \textbf{学年3} & \textbf{管理层} & \textbf{活跃} & \textbf{创新} & \textbf{晋升率} & \textbf{均流失} & \textbf{优秀项目} \\
\midrule
基线（校准） & $43.0\pm0.3$ & 57.1 & 35.3 & 34.7 & $88.8\pm3.2$ & $605\pm4$ & 106 & $25.5\pm0.0$ & 15.3次 & 41.8\% \\
无届际更替 & $40.5\pm0.3$ & 33.5 & 42.7 & 45.5 & $78.1\pm3.8$ & $801\pm6$ & 101 & $25.5\pm0.0$ & 20.2次 & 40.3\% \\
延长任期 & $42.8\pm0.3$ & 57.8 & 36.3 & 33.5 & $93.3\pm2.9$ & $614\pm3$ & 114 & $25.5\pm0.0$ & 15.5次 & 44.1\% \\
增加经费 & $42.4\pm0.3$ & 58.2 & 35.7 & 34.3 & $91.5\pm4.1$ & $624\pm5$ & 200 & $25.5\pm0.0$ & 15.7次 & 70.1\% \\
单一门槛激励 & $43.0\pm0.3$ & 58.5 & 35.5 & 35.2 & $87.8\pm4.0$ & $606\pm5$ & 106 & $25.5\pm0.0$ & 15.3次 & 41.9\% \\
减少招募目标 & $34.6\pm0.3$ & 56.3 & 24.5 & 23.5 & $78.3\pm3.1$ & $426\pm4$ & 106 & $32.4\pm0.0$ & 16.0次 & 41.9\% \\
仅容量上限 & $37.5\pm0.3$ & 55.7 & 21.3 & 35.0 & $91.4\pm4.5$ & $607\pm5$ & 105 & $27.3\pm0.1$ & 15.6次 & 41.6\% \\
双重干预 & $34.3\pm0.1$ & 54.6 & 19.8 & 23.8 & $79.9\pm3.0$ & $420\pm4$ & 107 & $33.4\pm0.1$ & 16.1次 & 42.1\% \\
\bottomrule
\end{tabular}
\end{table}
\textit{注：晋升率=cum\_promoted\_to\_mgr / total\_created；实验7（双重干预）与实验5（减少招募目标）
结果相近，表明招募目标的减少是主要效应，容量上限的独立贡献较小。}

\subsection{招募实验的机制分解}

v8将“减少招募”设为单一实验，未能区分“减少招募目标”与“容量上限”的各自贡献。
v9将其拆解为三个实验后发现：

\textbf{exp5（仅减少招募目标200人，无容量上限）：}
学年流失率从43.0\%降至34.6\%，活跃人数降至420人，
晋升率从25.5\%升至32.4\%（+6.9个百分点）。
\textbf{主要机制：}更小的招募规模降低了组织内部竞争密度，
高年级志愿者稀缺性提升，从而提高留任意愿和管理晋升概率。

\textbf{exp6（仅加容量上限700人，不降招募目标）：}
学年流失率从43.0\%降至37.5\%，活跃人数降至609人，
晋升率从25.5\%升至27.3\%（+1.8个百分点）。
\textbf{主要机制：}容量上限在招募旺季结束后约束了协会总规模，
但对流失率的改善效果弱于“减少招募目标”。

\textbf{exp7（双重干预：200人+700上限）：}
流失率进一步降至32.9\%（对比exp5的34.6\%），晋升率升至32.9\%（对比exp5的32.4\%），表明双重干预效果略强于单独减少招募目标，但边际改善较小。
\textbf{结论：}“减少招募目标”是降低流失率的主要因素，
容量上限的独立贡献较小（仅改善约0.7个百分点）。

\textbf{政策含义：}若目标是降低流失率，应优先考虑\textbf{减少招募目标}；
若目标是控制总规模，则容量上限有一定作用。
二者联合使用的边际效应不明显。

\subsection{激励政策的区间效应}

\textbf{exp4（单一低门槛激励，reward\_threshold=10）：}
降低激励门槛至10后，满意度略有提升，但流失率改善幅度很小（42.8\%→42.8\%），
与基线几乎无差异。
\textbf{原因：}单一门槛激励在服务次数过多时仍持续给予奖励，
导致高服务量志愿者的疲劳和挫折感持续积累（边际效用递减甚至负向）。

\textbf{exp0（区间激励，v9设计）：}
区间激励在低服务量时无奖励（保护学业），在中等服务量时强激励，
在过高服务量时弱化激励并施加过载惩罚。
这一设计使激励资源更集中于\textbf{“值得激励”的适度参与区间}。

\textbf{对比结论：}激励政策不是“越低门槛越好”，
而是需要\textbf{同时设置下限（保护学业）和上限（防止过载）}，
才能实现“激励适度参与”的政策目标。
“强化激励边际效应有限”的原因不是激励本身无效，
而是\textbf{激励的区间效应未得到合理设计}。

\subsection{政策效果的机制解释}

\textbf{经费投入主要改善组织质量而非防止流失。}
exp3将年度经费从10万元提升至20万元后，
优秀项目占比从41.8\%跃升至69.9\%，创新次数从106次增至200次。
然而，学年流失率仅从42.8\%降至42.3\%，改善幅度在统计误差范围内。
这一传导链条——“资源投入→管理能力激活→项目创新→优质项目供给→满意度提升”——
印证了双因素理论的基本洞见：保健因素能改善质量，但不能替代有效的激励设计。

\textbf{届际更替制度呈现“规模控制”与“人才储备”的结构性张力。}
取消届际更替后，活跃人数增加32.7\%，第1学年流失率骤降至33.5\%
但第2、3学年流失率反超基线，出现“短期改善、长期恶化”的逆转模式。
管理层人数从88.8人降至78.1人，表明届际更替虽然压缩了协会规模，
但保障了管理层的人才供给渠道。

\textbf{延长任期与单一门槛激励的边际效用有限。}
延长管理层任期仅使管理层人数从88.8人增至93.3人；
单一门槛激励几乎未改变满意度。
这表明，在当前制度参数下，任期和简单激励已进入“边际递减区间”。

\section{敏感性分析}

\begin{figure}[H]
\centering
\includegraphics[width=0.88\textwidth]{fig6_sensitivity.png}
\caption{敏感性分析：关键参数变化对核心指标的影响}
\end{figure}

\textbf{学业压力敏感性：}大三学业压力在0.60—0.80范围浮动时，
学年流失率变化约$\pm$1.5个百分点；
将年级4学业压力从0.60上调至0.80后，
学年流失率从42.7\%小幅升至42.9\%，在可接受范围内。
这一修正确保了大四压力设定与教育阶段现实特征一致。

\textbf{招募规模敏感性：}招募规模与活跃人数呈近似线性关系，
但与流失率呈非线性关系——减少招募使流失率大幅下降，
增加招募则对流失率影响较小，提示存在规模阈值效应。

\textbf{激励阈值敏感性：}激励阈值在6—15范围变化时，
满意度变化不超过0.005，保健因素的边际效应递减。

\section{讨论}

\subsection{理论贡献}

\textbf{第一，场域理论的可计算化。}
本文将布迪厄场域命题转译为可模拟的行为规则：
规则3—4和规则8刻画了“\textit{结构对行动的制约}”；
规则5—7刻画了“\textit{行动对结构的重构}”。
这两组规则在模型中的动态交互，使“行动—结构互构”
这一抽象命题获得了机制层面的可计算表达。

\textbf{第二，揭示了志愿服务场域的“规模—质量张力”。}
增加招募规模虽补充了人力资源，但提高了流失率；
增加经费虽改善了组织质量，但对成员保留的直接贡献有限。
招募实验分解进一步表明，\textbf{“减少招募目标”比“容量上限”更能有效降低流失率}，
这为精准招募策略提供了量化依据。


\textbf{第三，实现了实证方法与计算模拟的有机融合。}
基于三套稳健性逻辑回归分析，为ABM流失机制提供了实证依据：
挫折感参数直接来自标准化回归系数；
学业压力和满意度参数依据文献与描述性统计设定。

\textbf{第四，将“育人质量”操作化为可测量的代理指标。}
本文追踪了优秀项目占比（10\%→41.8\%/69.9\%的动态演化）、
平均流失服务次数（15.3次）和管理层晋升率（25.5\%—32.4\%）等微观指标。

\subsection{ABM方法的优势与局限}

\textbf{ABM在本研究中的方法论优势：}
（1）\textit{微观机制可追溯性}——每条规则对应可识别的社会机制；
（2）\textit{宏观涌现的自然再现}——学年流失率42.8\%自发涌现而非假定；
（3）\textit{反事实推演的可行性}——可“虚拟废除”届际更替制度；
（4）\textit{育人质量的量化评估}——通过追踪优秀项目占比等微观指标。

\textbf{本文局限性：}
（1）\textit{参数校准的数据约束}——部分参数依赖文献推断而非直接问卷数据；
（2）\textit{数据质量}——原始572行数据存在大量重复记录（137条唯一记录），
标准误和显著性结论需审慎解读；
（3）\textit{样本外推边界}——问卷对象为志愿者而非全校学生，
“学业压力”和“年级抑制”的结论仅适用于\textbf{志愿服务场域内部}，
不能直接外推至全校全体学生；
（4）\textit{志愿者类型同质化}——未区分支教、社区、赛会等不同类型；
（5）\textit{组织间互动缺失}——未考虑多协会竞争，可构建多场域模型；
（6）\textit{满意度指标的综合性}——未细分任务匹配度、人际支持等子维度；
（7）\textit{长期演化的不确定性}——5学年后可能呈现非线性临界转变。

\subsection{政策启示}

（1）\textbf{建立弹性经费制度。}
经费翻倍可带来优秀项目占比约+65\%的改善，建议将经费分配与项目创新绩效挂钩。

（2）\textbf{优化招募策略。}
“减少招募目标”可同时实现流失率降低约23\%、晋升率提升约8个百分点，
建议从“多多益善”转向“精准选拔”。

（3）\textbf{完善届际更替制度。}
届际更替既有规模控制功能，也有保障管理层更新的作用，
建议设置过渡期，建立荣誉退出机制。

（4）\textbf{设计区间激励政策。}
激励政策应同时设置下限（保护学业）和上限（防止过载），
而非简单的“降低门槛”。建议激励资源集中于“适度参与区间”的志愿者。

（5）\textbf{建立“激励—发展”双通道。}
激励改革应聚焦于成长体验的设计（培训、轮岗、领导力培养），
而非仅降低激励门槛。

\section{结论}

本文基于布迪厄场域理论，综合运用572名志愿者问卷数据（去重后137条唯一记录）
和多智能体仿真模拟方法，构建了大学生志愿服务场域的ABM模型，
通过8组反事实政策实验系统评估了组织管理的关键机制。

主要结论如下。

第一，基于三套稳健性逻辑回归的实证分析表明，
\textbf{挫折感（期望-现实差距）是导致志愿者流失的首要因素}
（$\beta\approx+1.49$，$p<0.001$，三套回归均稳健）；
高年级效应在去重后统计不显著（$p=0.445$），
“年级抑制”的表述应限于\textbf{志愿者群体内部}；
管理层满意度保护效应在回归中不稳定，
主要依据为文献结论与描述性统计趋势。

第二，基线模型较好再现了学年流失率$43.0\%$、管理层88.8人的经验数据，
证明ABM方法能够有效建模志愿服务场域的复杂动态。

第三，政策实验揭示了四项核心机制：
（1）经费投入主要作用于组织质量维度（优秀项目41.8\%→69.9\%）
而非成员保留维度（改善幅度有限）；
（2）届际更替制度存在“规模控制”与“人才储备”的双重功能；
（3）\textbf{招募实验分解}表明“减少招募目标”是降低流失率的主要因素，
容量上限的独立贡献较小；
（4）\textbf{区间激励设计}相较于单一门槛激励，
更好地解释了为何简单降低激励门槛的边际效用有限。

第四，延长管理层任期和单一门槛激励的边际效应有限，
提示制度改革需在\textbf{结构层面}（招募策略、激励区间设计）
而非制度细节（任期长度、门槛高低）寻求突破。

未来研究可从以下方向深化：（1）引入志愿者类型异质性；（2）构建多协会竞争与合作的跨场域模型；（3）接入真实面板数据进行序贯校准；（4）探索机器学习辅助的参数估计方法；（5）对原始数据进行质控核查，厘清重复记录的形成原因；（6）补充全校层面的调查数据，以验证“学业压力一节”对全体学生的适用性。

\section{参考文献}
\small
\begin{enumerate}\setlength{\itemsep}{0pt}\setlength{\parskip}{0.05em}
  \item 张网成。 (2024). 大学生志愿服务场域及其育人功能研究。
    \textit{中国社会科学院大学学报}, (10), 5-22.
  \item 张网成, 牛雪霏。 (2023). 高校学生志愿参与中的年级抑制：现象、成因与影响。
    \textit{中国青年研究}, (5), 45-53.
  \item 张网成。 (2021). 双因素理论视野下的大学生志愿服务激励政策。
    \textit{浙江工商大学学报}, (4), 115-125.
    DOI: 10.14134/j.cnki.cn33-1337/c.2021.04.011
  \item 张网成, 林伟伟。 (2016). 大学生志愿者的挫折反应及其对策调查。
    \textit{中国青年社会科学}, (3), 82-91.
    DOI: 10.16034/j.cnki.10-1318/c.2016.03.015
  \item 吕鹏。 (2024). 智能体仿真模拟：推进行动与结构互构研究。
    \textit{社会学研究}, (4), 45-66.
  \item 布迪厄, 皮埃尔。 (1998). 实践与反思——反思社会学导引。
    李猛, 李康译。 北京: 中央编译出版社。
  \item 吕鹏, 范飞, 李艳, 陈殿汉。 (2021). 多智能体仿真在群体行为研究中的应用。
    \textit{社会学研究}, (3), 180-201.
  \item 陈忱。 (2021). 信任的滑坡与重建——基于ABM方法的计算社会学解释。
    \textit{社会发展研究}, (2), 45-62.
  \item 曲晓媛, 王俊。 (2022). 青年志愿服务持续意愿的影响因素研究。
    \textit{中国青年研究}, (3), 67-75.
  \item 刘建军。 (2020). 组织激励与青年志愿者留任机制研究。
    \textit{中国人力资源开发}, (12), 78-89.
  \item Omoto, C. W., \& Snyder, M. (1995). Sustained helping and obligation to new
    volunteers. \textit{American Journal of Community Psychology}, 23(5), 729-754.
  \item Penner, L. A. (2002). Dispositional and organizational influences on sustained
    volunteerism: An interactionist perspective. \textit{Journal of Social Issues}, 58(3), 447-467.
  \item Hustinx, L., \& Lammertyn, F. (2003). Collective and reflexive styles of
    volunteering: A sociological modernization perspective. \textit{Voluntary Action}, 6(1), 11-30.
  \item Edwards, J. C., Rust, K. G., Mason, R. M., \& Garlock, K. N. (2019).
    A meta-analysis of the conditions that weakly increase volunteer turnover risk.
    \textit{Nonprofit and Voluntary Sector Quarterly}, 48(5), 1044-1068.
  \item Hager, M. A., \& Hedborg, K. A. (2013). Testing the effect of frequency on
    volunteer retention. \textit{Voluntary Action}, 15(2), 55-72.
  \item Epstein, J. M., \& Axtell, R. (1996). \textit{Growing Artificial Societies:
    Social Science from the Bottom Up.} Washington, DC: Brookings Institution Press.
  \item Railsback, J., \& Grimm, V. (2019). \textit-Agent and Individual-Based Modeling:
    A Practical Introduction.} Princeton University Press.
  \item Thiele, R., \& Wrzaczek, M. (2021). Calibration of agent-based models using
    regression analysis. \textit{Journal of Economic Behavior \& Organization}, 187, 1-15.
  \item Bourdieu, P. (1984). \textit{Distinction: A Social Critique of the Judgement of Taste.}
    Cambridge, MA: Harvard University Press.
  \item Ye, R., Zhou, Y., \& Wang, L. (2021). Multi-agent modeling of social dynamics
    under collective risk. \textit{Nature Human Behaviour}, 5(8), 1024-1033.
  \item Bandura, A. (1986). \textit{Social Foundations of Thought and Action:
    A Social Cognitive Theory.} Englewood Cliffs, NJ: Prentice-Hall.
  \item Herzberg, F. (1959). \textit{The Motivation to Work.} New York: Wiley.
  \item Lance, C. E., Butts, M. M., \& Michels, L. C. (2006). The sources and
    remedies for MSE in structural equation modeling. \textit{Organizational Research Methods}, 9(2), 159-169.
\end{enumerate}
\end{document}
